\documentclass{journal}

\usepackage{graphicx}
\graphicspath{{../../drills}}


\usepackage{amsmath}
\usepackage[nopatch]{microtype}
\usepackage{booktabs}
\usepackage{hyperref}

\usepackage[left=2.45cm,top=3cm,right=2.45cm,bottom=2.5cm,bindingoffset=0cm]{geometry}

\title{Intercepting Jammers}

\author{}

\begin{document}
\maketitle
\noindent 

This session is about catching, incipient jammers.
It will also cover a range of jammer control movements, including hit outs and run-backs.

\section*{Overview}

\subsection*{Pre-Session: Off Skates Warm-up 5-10 Minutes}
Lightly making sure that the joints still work.

Focus is {\bf not} on building strength, cardio or endurance, just on limbering up. Do 5 to 10 reps of each, on each side. 

\begin{itemize}
\item Toe taps + one leg light squats
\item Arm, shoulder and neck stretches
\item Opening and closing hip rotations
\item Knee raise to superman 
\item Cross chest knee raises
\item Butt kicks, high knees, side shuffles
\end{itemize}


\subsection*{Warm Up - 10 Minutes}

This is a juking and hitting session, so we modify the warm up accordingly.

\begin{itemize}
\item Sticky skating - Regular, Figure 8, low
\item Weaves rampage style. 
\item \hyperref[drill:footwork/weave_carve_weave]{Weave, Carve, Weave} drill.
\item Side Surfing 
\item 1 minute laps, 1 minute backwards, 1 minute anti-derby, 1 minute anti-derby backwards.  
\end{itemize}

Followed by any dynamic stretching needed.
For a contact warm-up we will do one-foot blocking.

We will then level up the one foot blocking with the engaging skater practicing a drive. 


\subsection*{Blocking Warmup Drills}
\begin{itemize}
\item \hyperref[drill:one_on_one/developmental/spinning_blocking]{Spin blocking}
\item Ugly Blocking
\item \hyperref[drill:one_on_one/zonal_catching]{Zonal catching}.
\item Offence leap frog \hyperref[drill:drives/leap_frog]{drive variant}.
\end{itemize}

\subsection*{Jammer interception - 45 minutes}
Focus on line coverage, double taps, and jammers starting disengaged.
Jammers should try to not facilitate.


We will begin with a set of jammer running drills, here the blocker is aiming to take the jammer out, then run them back.

Jammers are instructed to try to pass in a single motion - i.e their first move should not be a juke if they're in contact, but they can juke afterwards.

\begin{itemize}
\item Starting in contact with the jammer, on a line or in the centre.
\item Starting three feet from the jammer on either a line or in the centre.
\item Starting five feet from the jammer on either a line or in the centre.
\end{itemize}


Next we add a second or third blocker depending on the jammer.
The goal here is for the blockers to catch, then to drive out and cut their own friend. 


\subsection*{Moving Walls - 45 minutes}
This next session is a revision on different driving drills.
Each drill should be run twice, once to feel out the movement, and the second time with a small modification.

\begin{itemize}
\item \hyperref[drill:drives/hip]{Hip drive} starting engaged
\item Hip drive from one to two feet 
\item \hyperref[drill:drives/shoulder]{Shoulder drive} 
\item \hyperref[drill:drives/low_shoulder]{Low shoulder drive}- with a focus on shoulder to butt cheek driving where possible. 
\item Shoulder drive into \hyperref[drill:drives/chest]{Chest drive}
\item \hyperref[drill:drives/shelf]{Shelf drive} starting engaged.
\end{itemize}

Blockers should attempt to resist, especially using the one foot blocking from earlier.
Blockers should act as if they were a butt in a wall - hence they cannot simply avoid the hit without displacing themselves 

The offence should be seeking to move the blocker at least one lane, but ideally to take them off the track. 

\subsection*{360 Offence - 10 minutes}

Mixing up the drives now, start with an offence engaged in a three wall. 

\begin{itemize}
\item Run the \hyperref[drill:three_wall/360_offence]{360 offence drill} - 30 seconds to 1 minute per offence, then a rest.
\item Move onto \hyperref[drill:three_wall/skewering]{skewering} practice.
\end{itemize}

This is a higher intensity drill, so cycling skaters for a rest is a good idea.


\subsection*{Jammer and Offence - Remainder of session}

Jammer starts disengaged, either 5, 10 or 15 feet.
Blockers should seek to work in the cutting their own friends drill on runbacks. 


\begin{itemize}
\item Initially just a jammer against a three wall
\item Reduce blockers to two wall
\item Increase blockers to a three wall, add an offence.
\item Increase blockers to a four wall - fourth blocker may do a four wall defence, or run counter-offence.
\end{itemize}



\pagebreak


\section*{Drill Notes}

\subsection*{Drill: Weave, Carve, Weave}
\label{drill:footwork/weave_carve_weave}

Starting in the centre of the track with a runup.
The first motion is a weave towards either line, followed by a sharp carve back to the centre, followed by a corrective weave to continue the forwards direction of travel.  

This should resemble a possible jammer juke, but is mostly to stress the edges under speed. 

% TODO: Figure


\input{../../drills/non_contact/stopping/dynamic_take}
\subsection*{Drill: Derby Stops}
\label{drill:stopping/derby_stops}

Around the track, the goal is to approach a line, and perform a sharp turn around toe-stop without going out, then to accellerate back towards the centre before taking the other line.

Side shuffles and grape-vines can be incorporated, mimicing a jamming passing movement.



\input{../../drills/contact/solo_blocking/developmental/spinning}
\subsection*{Offence Leap Frog: Driving}
\label{drill:drives/leap_frog}

This is a driving variant of the leap frog.

Start with a skater in the centre of the track, rather than the line.

Perform a run-up into a drive on the skater, attempting to get them to the edge of the track.
Then reset in the centre, swapping such that the driven skater is now driving. 

Swap which skater is on the inside and outside.



\begin{minipage}{0.38\textwidth}
\begin{center}
\includegraphics[scale=1,trim={0 1cm 0 4cm}, clip]{img/one_on_one/zonal_blocking/position}

{\it The track may be divided using cones. The blocker is strongly encouraged to suck back past the front cones into the jammer's space.} 
\end{center}
\end{minipage}
\hfill
\hspace{0.5cm}
\begin{minipage}{0.6\textwidth}

\subsection*{Drill: Zonal Catching}
\label{drill:one_on_one/zonal_catching}

% Nova Drill

Divide the track into thirds, with a middle, inside and outside lane. 
Place cones on these dividing lines marks, spaced about 50cm apart.

This drill requires two skaters, one acting as a jammer, one as a blocker, however this is a blocker focused drill. 
The jammer should sequentially move backwards around the cones, then attempt to pass the blocker in the next lane. 
The blocker should move between lanes, tracking and catching the jammer.

The goal of this drill is the footwork and the catch rather than the jammer attempting to pass the blocker.

The drill may be run with any form of catch.

\end{minipage}




\input{../../drills/contact/solo_blocking/driving/hip}
\subsection*{Drill: Shoulder Drive}
\label{drill:drives/shoulder}

This drill requires pairs of skaters. 

Starting engaged, one skater is driving the other blocking, with both skaters starting in the middle of the track. 

The driving skater should attempt to use their shoulder to drive the other blocker.  
The skater being driven should attempt to hold their position (rather than simply skating away).


In this case, the shoulder can be used to partially lift the other blocker.
The idea is to set up in the armpit or ribs of the other blocker, and leverage from this position.
As friction requires that they're applying weight into the floor, this drive effectively negates the blockers ability to engage their edges with one leg.    


Practice in pairs, and rotate skaters where possible - driving is best practiced against a wide range of different opponents.
Also don't forget to practice driving to both the inside and outside.


% TODO: Figure

\subsection*{Drill: Low Shoulder Drive}
\label{drill:drives/shoulder_low}

This drill requires pairs of skaters. 

Starting engaged, one skater is driving the other blocking, with both skaters starting in the middle of the track. 

The driving skater should attempt to use their shoulder to drive the other blocker.  
The skater being driven should attempt to hold their position (rather than simply skating away).


In this case, the shoulder can be used to partially rotate the other blocker. 
The driving skater should set their shoulder up next to the butt cheek of the blocker. Using this as a fulcrum they can then apply pressure to rotate the blocker. 
Once the blocker is rotated, they've effectively had their feet turned, as in the hip drive, and can now be pushed more easily. 

For this secondary drive, try moving from the shoulder drive into a chest drive to corale the blocker off the track. 


Practice in pairs, and rotate skaters where possible - driving is best practiced against a wide range of different opponents.
Also don't forget to practice driving to both the inside and outside.



% TODO: Figure

\subsection*{Drill: Chest Drive}
\label{drill:drives/chest}

This drill requires pairs of skaters. 

Starting engaged, one skater is driving the other blocking, with both skaters starting in the middle of the track. 

The driving skater should attempt to use their chest to drive the other blocker.  
The skater being driven should attempt to hold their position (rather than simply skating away).

The idea of this drive is to push from a chest-on-chest position, using a combination of the direction of play, and the momentum from an initial, alternative drive to get the skater moving.     

This is best started from a shoulder drive, either low or regular, then ends in a chest drive. 
Keeping one leg down-track from the blocker prevents them from rotating in that direction, ideally getting them onto toe-stops further limits their ability to get out of the situation, which is helped by lifting them with the initial shoulder drive. 


Practice in pairs, and rotate skaters where possible - driving is best practiced against a wide range of different opponents.
Also don't forget to practice driving to both the inside and outside.


% TODO: Figure

\subsection*{Drill: Shelf Drive}
\label{drill:drives/shelf}

This drill requires pairs of skaters. 

Starting engaged, one skater is driving the other blocking, with both skaters starting in the middle of the track. 

The driving skater should attempt to shelf and drive the other blocker. 
The skater being driven should attempt to hold their position (rather than simply skating away).


This is a mix of the hip and shoulder drive.
Setting up with a leg placed laterally beneath the blocker, and a shoulder in the same position as the shoulder drive, the idea is to perform a lateral directly into the blocker. 
It is important to not engage the toe stop under the target blocker, as this will burn off the momentum of the drive.
Additionally it is important to start the engagement with either the hip, or the shoulder - starting with the thigh under the other blocker is a back block if any impact is achieved, and the purpose of the drive is to cause impact.


Using a `skateboard' lateral push, or a toe-stop push continue to drive the other blocker off the track.
The shoulder and thigh set up prevents the blocker from getting low - limiting the friction they can generate to resist the push.



Practice in pairs, and rotate skaters where possible - driving is best practiced against a wide range of different opponents.
Also don't forget to practice driving to both the inside and outside.



% TODO: Figure



\subsection*{Drill: 360 Offence}
\label{drill:three_wall/360_offence}

This drill requires four skaters, three blockers and an opposing offence. 


The original drill uses two to three different drives, I'll be updating this to four.

\begin{itemize}
    \item Hip drive using lateral pushes
    \item Shoulder to hip drive  
    \item Chest drive
    \item Shelf drive using lateral pushes 
\end{itemize}

The goal is to chain these drives together to dismember a wall.

The offence starts in the centre of the wall, where the jammer would normally be.

The offence should then attempt to drive one of the butts, switch and drive the brace then disengage and reset to drive the other butt towards the centre of the track before switching to again driving the brace. 
Then reset in the centre and repeat in the opposite direction.

Meanwhile the wall should seek to reform and counter the incoming offence.

% TODO: Fix link
%A video can be found \hyperref[https://www.youtube.com/watch?v=wZ6JcTYCBYw]{here}.

\begin{minipage}{0.48\textwidth}
\begin{center}
\includegraphics[scale=1,trim={0 1cm 0 4cm}, clip]{img/offence/360/1}


{\it Offence from within the seam, pushing out then taking the brace.}
\end{center}
\end{minipage}
\begin{minipage}{0.48\textwidth}
\begin{center}
\includegraphics[scale=1,trim={0 1cm 0 4cm}, clip]{img/offence/360/2}


{\it Offence starting from the line, optionally skewering, then taking the brace.} 
\end{center}
\end{minipage}



{\it Progressions}

As an offence progression - drive both butts together off the track as part of the final movement as in the skewering drill.

%As a natural progression, add a jammer.
% TODO: Figure

\begin{minipage}{0.38\textwidth}
\begin{center}
\includegraphics[scale=0.6,trim={0 0 0 4cm}, clip]{img/offence/skewering/position}
{\it Jammer left with only the brace to challenge.} 
\end{center}
\end{minipage}
\hfill
\hspace{0.5cm}
\begin{minipage}{0.6\textwidth}
\subsection*{Drill: Skewering}
\label{drill:three_wall/skewering}

This drill requires four skaters, three blockers and an opposing offence. 

The blockers begin set up as a three-wall.
The offence's goal is to hit one of the butts into the other, then to launch into the brace. 

\begin{itemize}
\item Queue up from the inside or the outside of the track, lining the blockers up.
\item The initial hit into drive is to impart momentum onto the first blocker, allowing their own force to drive them into the second.
\item Once that blocker is no longer making contact with the offence, the next move is to take the brace.
\end{itemize}

The initial drive may be any of the hip, shoulder, chest or shelved drives.
When hitting the brace consider the advantages of either covering the same line as the initial approach, or setting up a threading situation. 


% TODO: Fix link
%A video can be found towards the end of \hyperref[https://www.youtube.com/watch?v=wZ6JcTYCBYw]{here}.
\end{minipage}




\end{document}
